% Main-text reduction; detailed sources and Green transposes are in the technical appendix.
\section{Reduced representation of the complete cubic response}\label{sec:stationary}
The complete map can be evaluated without an independent third-order field solve or a
second-order photon ray. Its three ingredients are a measurement on first-order fields,
a measurement Hessian on the known second-order fields, and a linear response to the
complete cubic source. The Einstein constraints make the last response acoustic and
wave-like: the hyperbolic scalar Green functions cancel before integration.

\paragraph{Angular and time conventions.}
Write $e=\bar t_\star$, $o=\bar t_o$, $\iota=\bar t_I$, with $0<\iota<e<o$, and
$X_0=n(o-t)$. For a quadrupole covector $\mathsf C$ (symmetric and trace free), define
\begin{equation}\label{eq:adjoint-weight-convention}
\begin{gathered}
 w=\frac{15}{2}n^T\mathsf Cn,\qquad W=\frac{w}{4\pi},\qquad
 J=\sky{w\Theta}=\int W\Theta\,\dd\Omega,\\
 C_K(t)=\int\dd\Omega\,W\,\TF(nn)e^{iK\cdot n(o-t)}.
\end{gathered}
\end{equation}
Thus $\sky{w\,\TF(nn)}=\int W\,\TF(nn)\dd\Omega=\mathsf C$.
Every full-solid-angle objective, source pairing and multiplier below and in
App.~\ref{app:cubic-technical} uses $W$; sky means use $w$. For other multipoles,
$W$ is the corresponding zero-mean harmonic weight. The notation $J_{abc}$ denotes
the polarised cubic response of this objective, not a power spectrum.

\subsection{The combined response}\label{sec:combined}
\begin{proposition}[Combined cubic response]\label{prop:combined}
For an arbitrary triple with $P=k_a+k_b+k_c\neq0$ and any zero-mean harmonic weight,
\begin{equation}\label{eq:combined}
\boxed{J_{abc}=J_{111,abc}-\sum_{a|bc}S_{{\rm meas},a}[\nu_{bc},A_{bc}]+J_{\rm lin}[y_{abc}].}
\end{equation}
Here $J_{111}$ is the complete measurement on the first-field-only family
$(\zeta,\nu,H)=\eps(g,a,h)$, including its physical roots, clock, frame and normalisation.
The three partitions $a|bc$ each have coefficient one. Their fields
$y_{bc}=(\nu_{bc},A_{bc})$ are the closed-form physical second-order fields of
\S\ref{sec:fields}, with the exact $K=0$ kernel for a zero pair momentum.
The functional $S_{\rm meas}$ is minus the reduced measurement Hessian; its explicit
no-probe-ray form is \eqref{eq:measurement-source}. Finally,
\begin{equation}\label{eq:combined-linear}
\begin{split}
 J_{\rm lin}[y_{abc}]={}&-3\int_e^o\ell_P(t)\Big(\nu_{abc}'-\frac{\nu_{abc}}t\Big)\dd t
 -i\int_e^o\mathcal P_P^T[C_P(t)]\cdot c^{(3)}_M(t)\,\dd t\\
 &-\frac12\int_e^o C_P:h^{{\rm TT}\prime}_{abc}\,\dd t,
\end{split}
\end{equation}
where $\ell_P=\hat P\cdot C_P\cdot\hat P$,
$\mathcal P_P^T[C]=2CP/P^2-(P\cdot C\cdot P)P/(2P^4)$, and the responses mean the
explicit source integrals
\[
 \nu_{abc}(t)=c_{abc}v_P(t)-\int_0^tG_N^+(t,s)c_N^{(3)}(s)\dd s,\qquad
 h^{\rm TT}_{abc}(t)=-2\int_0^t\frac{s}{t}\frac{\sin P(t-s)}P\,
 \mathrm{TT}[c_T^{(3)}](s)\dd s.
\]
Here $v_P=9u_1(Pt/\sqrt3)/P^2$, and $G_N^+$ is the retarded inverse of the lapse
operator $L_N$, with the sign opposite to the source kernel in \S\ref{sec:fields}.
The complete sources $c^{(3)}$ are local in the first- and second-order fields.
No independent third-order field solve, first field adjoint or second ray is needed.
\end{proposition}
\begin{proof}
Apply the third derivative chain rule to the exact field-reduced measurement:
$J_{abc}=J_{111,abc}+\sum_{a|bc}D_a(J_{{\rm red},y})[y_{bc}]+J_y[y_{abc}]$,
where $D_a$ differentiates along the frozen first direction $(y_a,g_a)$.
Each middle term is the measurement Hessian on that direction and the second-order
state; by definition it is $-S_{{\rm meas},a}[y_{bc}]$.
The cubic field equation is
$\mathcal L_Py_{abc}+c^{(3)}=0$, with its constrained initial row.
The source is compatible by Prop.~\ref{prop:compat}. Therefore
Prop.~\ref{prop:constraint-response} replaces its non-TT spatial response by the
momentum integral. Contracting its time derivative with $-C_P/2$ uses
$\mathcal P_P^T[C_P]\cdot P=3\ell_P/2$ and gives \eqref{eq:combined-linear}.
Equivalently, the four background-adjoint blocks of \eqref{eq:stationary-field}
combine into this single pairing with the complete cubic source.
\end{proof}

This supplies the cubic ingredient of \eqref{eq:primary}:
$R_\ell=\frac1{8\pi i^\ell}\int_{-1}^1J_{abc}\,\dd\mu$, with weight $W=P_\ell(\hat z\cdot n)$ and
triple $(k_0\hat z,k_0\hat q_\mu,-k_0\hat q_\mu)$.

The direct measurement $J_{111}$ is given without a second ray in
\eqref{eq:direct-U3}--\eqref{eq:second-clock}. Its frame term is already included and must
not be added again. Counting independent times after distributing products, every term of
\eqref{eq:combined} has at most two time integrations except $J_{111}$, whose ray products
have three. On the shell, the top-order response reduces to the acoustic nine-moment
form \eqref{eq:adjoint-nine-moment}. These are exact reductions, not evaluations of the
surviving time integrals. At total momentum $P=0$, replace the projected response by the
regular Cartesian $K=0$ integral of \S\ref{sec:fields}; no hyperbolic kernel is needed there
either. Implementation comparisons are collected in App.~\ref{app:validation}.

\subsection{Why the constraints remove hyperbolic kernels}\label{sec:constraint-rows}
The Cartesian spatial operator has eigenvalue $-K^2/3$ on $\TF(KK)$.
Its unconstrained scalar Green function is hyperbolic. Physical sources do not excite an
independent growing scalar mode: the momentum rows replace the scalar/vector evolution
rows, while the initial Hamiltonian row fixes the regular lapse coefficient.

\begin{proposition}[Source compatibility]\label{prop:compat}
Let $(c_N,c_M,c_T,c_{\rm Ham})$ be the order-$r$ sources of \S\ref{sec:fields}, with
$L_{\rm Ham}=4K^2\nu+12\nu'/t+12\nu/t^2-K\cdot A\cdot K$.
If all lower orders solve the field equations, then
\begin{align}
 c_{\rm Ham}'-\tfrac4tc_N&=2iK\cdot c_M,\label{eq:compat-1}\\
 c_M'+\tfrac2tc_M&=iK\cdot c_T-\tfrac{iK}6(c_{\rm Ham}+4c_N).\label{eq:compat-2}
\end{align}
\end{proposition}
\begin{proof}
Matter conservation is built in (Lemma~\ref{lem:frame}), so the contracted Bianchi
identity holds exactly. At order $r$, lower residuals vanish and its action is the
background Bianchi operator on $\mathcal L_K[y_r]+c^{(r)}$.
For arbitrary off-shell $\nu,A$, direct differentiation of the linear rows gives
\[
 (L_{\rm Ham})'-\frac4tL_N-2iK\cdot L_M=0,\qquad
 L_M'+\frac2tL_M-iK\cdot L_T+\frac{iK}{6}(L_{\rm Ham}+4L_N)=0.
\]
Subtracting these identities leaves \eqref{eq:compat-1}--\eqref{eq:compat-2} for the sources.
The lower-order on-shell assumption is essential; an arbitrary mixed residual need not
satisfy these source identities on its own.
\end{proof}

\begin{proposition}[Constraint-row response]\label{prop:constraint-response}
For $K\neq0$ and compatible sources, the regular solution satisfies
\begin{equation}\label{eq:constraint-response}
 \boxed{K\cdot A(t)=V(t)=\int_0^t\big[4K(\nu'-\nu/s)+2ic_M\big]\dd s,\qquad
 A-\mathrm{TT}[A]=\mathcal P_K[V],}
\end{equation}
where
$\mathcal P_K[V]=\TF(K\otimes V+V\otimes K)/K^2-(K\cdot V)\TF(KK)/(2K^4)$.
The lapse $\nu$ is the acoustic response with its regular coefficient fixed by the initial
Hamiltonian row. Equivalently,
$h_S=6\nu-6\int_0^t\nu/s\,\dd s+(3i/K^2)\int_0^tK\cdot c_M\,\dd s$.
Only the acoustic kernel, the TT wave kernel and one cumulative integral occur.
\end{proposition}
\begin{proof}
$\mathcal P_K[V]$ is trace free, has no TT part, and $K\cdot\mathcal P_K[V]=V$.
Contract the Cartesian equation with $K$ and use
$K\cdot\mathfrak B_K\mathcal P_K[V]=-(K\cdot V)K/3$.
The transverse equation is \eqref{eq:compat-2} transverse to $K$.
For the longitudinal part put $E=12(\nu'/t+\nu/t^2)+c_{\rm Ham}$.
The lapse equation and \eqref{eq:compat-1} give
$E'=-4K^2\nu/t+2iK\cdot c_M$; the initial Hamiltonian row gives
$E(0)=36c_r+c_{\rm Ham}(0)=0$. The remaining terms are the longitudinal part of
\eqref{eq:compat-2}. Thus $\mathcal P_K[V]$ solves the non-TT Cartesian equation with
zero regular data, and uniqueness gives the result.
\end{proof}
The inverse powers of $K$ in $\mathcal P_K$ are removable only in the recombined physical
response. At $K=0$ use the Cartesian kernel $s(t-s)/t$, which is already non-hyperbolic.
On the shell this construction gives the single acoustic integral for $F_S$ in
Prop.~\ref{prop:shell-acoustic}. The weak transpose retains explicit emission and observer
point terms (\eqref{eq:constraint-transpose}); they cannot be replaced by zero terminal data.

\subsection{The stationary explanation}
Write the independent residual system of the complete experiment as $F(y,g)=0$.
The state $y$ includes fields, rays, roots, clock, frame and normalisation. Let
$A=F_y|_0$ have the regular inverse fixed by the physical boundary conditions and set
\[
 L=J+\langle\lambda,F\rangle,\quad A^*\lambda_0=-J_y|_0,\quad
 Ay_a=-F_g[g_a],\quad A^*\lambda_a=-D_a(L_{0,y}),
\]
where $D_a$ differentiates along the frozen first direction $(y_a,g_a)$ and the pairing
includes every boundary residual.
\begin{proposition}[Polarised stationary identity]\label{prop:stationary}
\begin{equation}\label{eq:stationary-polarised}
 \boxed{J_{abc}=D_aD_bD_cL_0+\sum_{a|bc}\langle\lambda_a,D_bD_cF\rangle.}
\end{equation}
No second or third state response and no second adjoint appears. Each partition has
coefficient one; for a repeated direction, $J'''=D^3L_0+3\langle\lambda_1,D^2F\rangle$.
\end{proposition}
\begin{proof}
Differentiate along the implicit solution. The background adjoint removes $y_{abc}$.
The remaining second-state terms satisfy
$\langle D_aL_{0,y},y_{bc}\rangle=-\langle\lambda_a,Ay_{bc}\rangle
=\langle\lambda_a,D_bD_cF\rangle$.
\end{proof}
This identity trades second-order fields for first adjoints. Formula \eqref{eq:combined}
trades back to the already closed-form second-order fields and combines the background
blocks before applying compatibility. The two representations serve different purposes:
the stationary form exposes boundary duality, while the combined form is the smaller
physical source representation. Their complete sources and boundary covectors are in
App.~\ref{app:cubic-technical}.

\subsection{Evaluated consequences}\label{sec:stationary-evaluated}
\paragraph{The cubic quadrupole through four gradients.}
Scale shell momenta by $\kappa$, align the external unit vector with $(1,0,0)$ and choose
$\mathsf C=\hat k\hat k-\mathbf1/3$. The complete internally averaged cubic kernel is
\begin{equation}\label{eq:stationary-four-gradient}
 \overline J_{abc}=\kappa^2\frac{4(o^2-e^2)}9
 +\frac{\kappa^4}{3402}\big[3984e^4-6256e^3o+4692e^2o^2-1853o^4-567\iota^4\big]+O(\kappa^6).
\end{equation}
The homogeneous-lapse partition contributes
$-\kappa^4(3o^4+5\iota^4-8e^4)/54$ and does not cancel. By contrast the homogeneous
tensor drops out of the entire isotropic contraction: its response is proportional to
$Q_{\hat q}=\hat q\hat q-\mathbf1/3$, and the remaining linear functional is independent
of $\hat q$, so $\ens{Q_{\hat q}}_{\hat q}=0$. The uncontracted map retains that tensor.
The uncontracted four-gradient result agrees as a polynomial identity with the curvature
map of \S\ref{sec:next-gradient}; both gradient orderings are needed in $C^{(13)}$.

\paragraph{The exact dilation channel.}
For one finite momentum and two exactly zero momenta, $Y=e^sX$ maps
$g=s+\eps e^{ip\cdot X}$ to $\eps e^{ie^{-s}p\cdot Y}$ and preserves the whole experiment,
including its constrained seed, physical clock and frame. Uniqueness of the regular branch gives
\begin{equation}\label{eq:dilation-channel}
 \boxed{\Theta^{(3)}(p,0,0;n)=(p\cdot\partial_p)^2\Theta^{(1)}(p;n),\qquad
 T^{(3)}_\ell(k,0,0)=(k\partial_k)^2T_\ell(k).}
\end{equation}
The derivative acts on the transfer profiles and the ray phase. By Prop.~\ref{prop:Tl} the
right side is explicit. This exact zero-momentum channel is not the nonzero-radius shell
contraction and does not authorize exchanges of soft limits with spectral integrals.

\paragraph{The first nonzero observer-centred cubic coefficient.}
With the same quadrupole weight, isotropy removes the $(q,-q)$ and single-$q$ dilation terms.
Thus $\ens{\mathcal K^c_3}=\overline J_{abc}+2\mathfrak D\ens{\mathcal K_2(k,q)}
+\mathfrak D^2\mathcal K_1(k)$. Substitution of the four-gradient quadratic map and the
linear transfer cancels the degree-two term identically and gives
\begin{equation}\label{eq:centred-kappa4}
\boxed{\ens{\mathcal K^c_3(k,q,-q)}_{\hat q}=\frac{k^4}{270}\Big[4(e-o)(12e^3-8e^2o+7eo^2+7o^3)
+5(o^4-\iota^4)\Big]+O(k^6).}
\end{equation}
The first part vanishes with the photon path; the second is a pure observer-clock term.
Its contribution to $\frac13C^{(13)}_2$ at degree six is
$(4\pi/25)\sigma^4(o^2-e^2)k_0^2\ens{\mathcal K^c_3}/9$.

\begin{remark}[Initial-event dependence at fixed measurement labels]\label{rem:iota}
The clock roots are $t_Ie^{\nu(t_I,0)}=\bar t_I$ and
$\int_{t_I}^{t_o}te^{\nu(t,0)}\dd t=(\bar t_o^2-\bar t_I^2)/2$.
Holding $\bar t_o$ fixed while varying $\bar t_I$ (thus adjusting the prescribed elapsed
proper time) gives exactly
\[
 \left.\frac{\partial t_o}{\partial\bar t_I}\right|_{\bar t_o}
 =\frac{\bar t_I}{t_oe^{\nu_o}}
 \left[\frac{e^{-\nu_I}}{1+t_I\nu'_I}-1\right].
\]
At fixed elapsed proper time $\tau_o$, the same formula holds without the $-1$.
For the symmetric-frame coordinate map, initial-event dependence is entirely clock composition:
$\partial_{\bar t_I}\Theta_{\rm sym}
=(\partial_{t_o}\Theta^{\rm coord}_{\rm sym})\,\partial_{\bar t_I}t_o$.
For the specified Fermi--Walker map on the shell this identity holds through cubic amplitude,
with an $O(\eps^4)$ remainder, since $[H_1,H_1']=0$ there. Mixed higher-order commutators need
not vanish. Nevertheless the two frames differ by a global sky rotation at each observation
event, so every realization's $\sum_m|a_{\ell m}|^2$ is exactly unchanged. Frame independence
of the spectrum does not require all-orders equality of the maps.
At fixed $\bar t_o$, the first-order clock coefficient obeys
$\partial_{\bar t_I}D_1=-\bar t_I[\nu_1(\bar t_I,0)+\bar t_I\nu'_1(\bar t_I,0)]/\bar t_o$.
This accounts for the $\iota^4$ dependence in \eqref{eq:centred-kappa4} and the degree-six
cross term, localising it in the clock and initial-event sector.
\end{remark}
